\section{Conclusões}
\label{sec:conclusoes}

Este trabalho investigou a transferência de aprendizado com EfficientNet-B0
para a classificação de alimentos no Food-101, mantendo os pesos das
convoluções congelados. O treinamento do
classificador em conjunto com os parâmetros de escala e deslocamento do
Batch Normalization, combinado com Stochastic Depth, apresentou o melhor
desempenho: 79,43\% de acurácia top-1 e 94,86\% de acurácia top-5 no
teste. O ganho em top-1 sobre o modelo de referência foi de 6,65 pontos
percentuais. O treinamento dos 42.016 parâmetros adicionais do Batch
Normalization, sem Stochastic Depth, já havia proporcionado um ganho de
6,15 pontos percentuais; a ativação dessa regularização acrescentou
0,50 ponto percentual.

A atualização apenas das médias e variâncias acumuladas apresentou
desempenho inferior ao modelo de referência. Nas condições avaliadas, os
resultados favorecem a adaptação conjunta das estatísticas e dos parâmetros
do Batch Normalization, em vez da atualização exclusiva das estatísticas.
Essa observação se limita ao protocolo adotado, com 10 épocas e uma
semente por configuração, sem permitir estabelecer a estabilidade do ganho
adicional proporcionado pelo Stochastic Depth.

Como continuidade, propõe-se repetir os experimentos com diferentes
sementes, avaliar o treinamento por mais épocas e investigar o
descongelamento parcial das convoluções. Essas comparações poderão
esclarecer a estabilidade dos ganhos observados e o potencial de uma
adaptação mais ampla do modelo ao reconhecimento de alimentos.
